\chapter{\LUA\ 回调}\label{callbacks}

\section{注册回调}

\topicindex{callbacks}

\libindex{register}
\libindex{list}
\libindex{find}

本库包含用于注册、查找和列出回调的函数。回调是在明确定义的位置被调用的
\LUA\ 函数。回调分为两类：一类与现有功能混合，另一类（启用后）替换功能。
在大多数情况下，第二类应与内置功能表现相似，因为下一步需要特定的数据。
例如，你可以替换断字例程。该函数接收一个可以（或不可以）断字的列表。
最终列表必须有效，并且（通常）用于构造段落。另一个函数可以替换连字构建器
和|/|或字距调整器。做其他事情也是可能的，但最终可能不会给用户带来预期的结果。

你需要做的第一件事是注册一个回调：

\begin{typing}
id, error =
    callback.register(<string> callback_name, <function> func)
id, error =
    callback.register(<string> callback_name, nil)
id, error =
    callback.register(<string> callback_name, false)
\end{typing}

这里的 \texttt{callback\_name} 是预定义的回调名称，见下文。该函数返回回调的
内部 \type {id}；如果回调无法注册，则返回 \type {nil}。在后一种情况下，
\type {error} 包含错误消息，否则为 \type {nil}。

\LUATEX\ 会将回调函数内部化，因此即使你意外地重新定义了一个函数也没有关系。

回调赋值始终是全局的。你可以使用特殊值 \type {nil} 代替函数来清除回调。

为了获得少量速度提升，你可以将布尔值 \type {false} 赋给与文件无关的回调，
这样做将阻止 \LUATEX\ 执行其默认情况下会执行的任何操作（当完全没有注册
回调函数时）。请注意：除非你 \notabene {确切地}知道自己在做什么，
否则这可能会导致各种问题！

\begin{typing}
<table> info =
    callback.list()
\end{typing}

表中的键是已知的回调名称，值是一个布尔值，其中 \type {true} 表示该回调
当前已设置（处于活动状态）。

\begin{typing}
<function> f = callback.find(callback_name)
\end{typing}

如果回调未设置，\type {find} 返回 \type {nil}。

\section{文件发现回调}

本小节中记录的行为被认为是稳定的，即不会再有向后不兼容的更改。

\subsection{\cbk {find_read_file} 与 \cbk {find_write_file}}

你的回调函数应遵循以下约定：

\begin{typing}
<string> actual_name =
    function (<number> id_number, <string> asked_name)
\end{typing}

参数：

\begin{itemize}
\item[\\texttt{id\_number}]

对于日志文件或 \prm {input} 文件，此数字为零。对于 \TEX 的 \prm {read}
或 \prm {write}，该数字加一，因此 \type {\read0} 变为~1。

\item[\\texttt{asked\_name}]

这是用户提供的文件名，由 \prm {input}、\prm {openin} 或 \prm {openout}
找到。

\end{itemize}

返回值：

\begin{itemize}
\item[\\texttt{actual\_name}]

这是实际使用的文件名。对于 \TEX 读入的第一个文件，你必须确保返回的
\type {actual_name} 带有扩展名，并且适合用作 \type {jobname}。否则，
你将不得不在 \LUATEX\ 结束后手动修复日志文件和输出文件的名称，并且
最终的格式文件名也会变得混乱。这是因为这些文件名依赖于 jobname。

如果找不到文件，你必须返回 \type {nil}。

\end{itemize}

\subsection{\cbk {find_font_file}}

\topicindex{callbacks+font files}

你的回调函数应遵循以下约定：

\begin{typing}
<string> actual_name =
    function (<string> asked_name)
\end{typing}

\type {asked_name} 是一个 \OTF\ 或 \TFM\ 字体度量文件。

如果找不到文件，返回 \type {nil}。

\subsection{\cbk {find_output_file}}

\topicindex{callbacks+output file}

你的回调函数应遵循以下约定：

\begin{typing}
<string> actual_name =
    function (<string> asked_name)
\end{typing}

\type {asked_name} 是用于写入的 \PDF\ 或 \DVI\ 文件。

\subsection{\cbk {find_format_file}}

\topicindex{callbacks+format file}

你的回调函数应遵循以下约定：

\begin{typing}
<string> actual_name =
    function (<string> asked_name)
\end{typing}

\type {asked_name} 是用于读取的格式文件（用于写入的格式文件始终在当前
目录中打开）。

\subsection{\cbk {find_vf_file}}

\topicindex{callbacks+font files}

类似于 \cbk {find_font_file}，但用于虚拟字体。这同时适用于 \ALEPH 的
\OVF\ 文件和传统的 Knuth 式 \VF\ 文件。

\subsection{\cbk {find_map_file}}

\topicindex{callbacks+font files}

类似于 \cbk {find_font_file}，但用于 map 文件。

\subsection{\cbk {find_enc_file}}

\topicindex{callbacks+font files}

类似于 \cbk {find_font_file}，但用于 enc 文件。

\subsection{\cbk {find_pk_file}}

\topicindex{callbacks+font files}

类似于 \cbk {find_font_file}，但用于 pk 位图文件。此回调接受两个参数：
\type {name} 和 \type {dpi}。在你的回调中，你可以决定查找：

\begin{typing}
<base res>dpi/<fontname>.<actual res>pk
\end{typing}

但其他策略也是可行的。由你来找到一个与该规格匹配的 ``合理
位图文件。

\subsection{\cbk {find_data_file}}

\topicindex{callbacks+data files}

类似于 \cbk {find_font_file}，但用于嵌入文件（\type {\pdfobj file '...'}）。

\subsection{\cbk {find_opentype_file}}

\topicindex{callbacks+font files}

类似于 \cbk {find_font_file}，但用于 \OPENTYPE\ 字体文件。

\subsection{\cbk {find_truetype_file} 与 \cbk {find_type1_file}}

\topicindex{callbacks+font files}

你的回调函数应遵循以下约定：

\begin{typing}
<string> actual_name =
    function (<string> asked_name)
\end{typing}

\type {asked_name} 是一个字体文件。此回调在 \LUATEX\ 构建其所需字体文件
的内部列表时被调用，因此实际的调用时机可能会让你感到意外。你的返回值
稍后会被反馈到匹配的 \type {read_file} 回调中。

奇怪的是，\cbk {find_type1_file} 也用于 \OPENTYPE{}（\type {OTF}）字体。

\subsection{\cbk {find_image_file}}

\topicindex{callbacks+image files}

你的回调函数应遵循以下约定：

\begin{typing}
<string> actual_name =
    function (<string> asked_name)
\end{typing}

\type {asked_name} 是一个图像文件。你的返回值用于从硬盘打开文件，
因此请确保返回的内容是你的操作系统认为的有效文件名。

\section{文件读取回调}\label{iocallback}

本小节中记录的行为被认为是稳定的，即不会再有向后不兼容的更改。

\subsection{\cbk {open_read_file}}

\topicindex{callbacks+opening files}

你的回调函数应遵循以下约定：

\begin{typing}
<table> env =
    function (<string> file_name)
\end{typing}

参数：

\begin{itemize}
\item[\\texttt{file\_name}]

由前一个 \cbk {find_read_file} 返回的文件名；如果没有定义此类回调，
则为 \type {kpse.find_file()} 的返回值。

\end{itemize}

返回值：

\begin{itemize}
\item[\texttt{env}]

这是一个表，其中至少包含一个必需的和一个可选的回调函数，用于此文件。
必需的字段是 \type {reader}，关联的函数将在每读取一行新内容时被调用一次；
可选的字段是 \type {close}，当 \LUATEX\ 处理完该文件时将被调用一次。
\LUATEX\ 从不查看表的其余部分，因此你可以用它来存储你私有的逐文件数据。
这两个回调函数都将该表作为它们唯一的参数。

\end{itemize}

\subsubsection{\type {reader}}

\topicindex{callbacks+reader}

每当 \LUATEX\ 需要从文件中读取新的输入行时，就会运行此函数。

\begin{typing}
function(<table> env)
    return <string> line
end
\end{typing}

你的函数应返回一个字符串或 \type {nil}。值 \type {nil} 表示已到达文件
末尾，并将使 \TEX\ 接下来调用可选的 \type {close} 函数。

\subsubsection{\type {close}}

\topicindex{callbacks+closing files}

当 \LUATEX\ 决定关闭文件时，将运行此可选函数。

\begin{typing}
function(<table> env)
end
\end{typing}

你的函数不应返回任何值。

\subsection{通用文件读取器}

\topicindex{callbacks+readers}

有一组回调用于加载二进制数据文件。它们都使用相同的接口：

\begin{typing}
function(<string> name)
    return <boolean> success, <string> data, <number> data_size
end
\end{typing}

\type {name} 通常是完整路径名，由文件发现回调之一或 \type
{kpse.find_file()} 的内部版本返回。

\begin{itemize}
\item[\texttt{success}]

当发生致命错误时返回 \type {false}（例如，当文件终究找不到时）。

\item[\texttt{data}]

构成该文件的字节。

\item[\\texttt{data\_size}]

\type {data} 的长度，以字节为单位。

\end{itemize}

如果找到了文件但读取时出现问题，则返回空字符串和零。

函数列表如下：

\begin{longtable}{l p{5cm}}
\toprule 函数  & 用途 \\
\midrule
\type{read_font_file}     & ofm 或 tfm 文件 \\
\type{read_vf_file}       & 虚拟字体 \\
\type{read_map_file}      & map 文件 \\
\type{read_enc_file}      & 编码文件 \\
\type{read_pk_file}       & pk 位图文件 \\
\type{read_data_file}     & 嵌入文件（如 \PDF\ 对象所支持的） \\
\type{read_truetype_file} & \TRUETYPE\ 字体文件 \\
\type{read_type1_file}    & \TYPEONE\ 字体文件 \\
\type{read_opentype_file} & \OPENTYPE\ 字体文件 \\
\bottomrule
\end{longtable}

\section{数据处理回调}

\subsection{\cbk {process_input_buffer}}

\topicindex{callbacks+input buffer}

此回调允许你在 \LUATEX\ 实际开始查看行输入缓冲区之前更改其内容。

\begin{typing}
function(<string> buffer)
    return <string> adjusted_buffer
end
\end{typing}

如果你返回 \type {nil}，\LUATEX\ 将假装你的回调从未发生。你可以从中
获得少量处理时间。此回调不替换任何内部代码。

\subsection{\cbk {process_output_buffer}}

\topicindex{callbacks+output buffer}

此回调允许你在 \LUATEX\ 因 \prm {write} 命令而实际开始将行输出缓冲区
写入文件之前更改其内容。它仅在输出到实际文件时被调用（即不包括日志、
终端以及所谓的 \prm {write} 18 调用）。

\begin{typing}
function(<string> buffer)
    return <string> adjusted_buffer
end
\end{typing}

如果你返回 \type {nil}，\LUATEX\ 将假装你的回调从未发生。你可以从中
获得少量处理时间。此回调不替换任何内部代码。

\subsection{\cbk {process_jobname}}

\topicindex{callbacks+jobname}

此回调允许你更改 \TEX 中由 \prm {jobname} 给出、以及 Lua 中由
\type {tex.jobname} 给出的作业名。它不影响内部作业名，也不影响输出文件
或日志文件的名称。

\begin{typing}
function(<string> jobname)
    return <string> adjusted_jobname
end
\end{typing}

唯一的参数是实际的作业名；你不应在此函数内使用 \type {tex.jobname}，
否则可能发生无限递归。如果你返回 \type {nil}，\LUATEX\ 将假装你的回调
从未发生。此回调不替换任何内部代码。

\section{节点列表处理回调}

节点和节点列表的描述见~\ref{nodes}。

\subsection{\cbk {contribute_filter}}

\topicindex{callbacks+contributions}

当 \LUATEX\ 向列表添加内容时，会调用此回调：

\begin{typing}
function(<string> extrainfo)
end
\end{typing}

该字符串报告组代码。由此你可以推断出你可以处理哪个列表。

\begin{longtable}{l p{5cm}}
\toprule 值             & 说明                                  \\
\midrule
\type{pre_box}    & 正在添加行间材料            \\
\type{pre_adjust} & 正在添加 \prm {vadjust} 材料 \\
\type{box}        & 正在添加一个排版好的盒子（始终被调用） \\
\type{adjust}     & 正在添加 \prm {vadjust} 材料 \\
\bottomrule
\end{longtable}

\subsection{\cbk {buildpage_filter}}

\topicindex{callbacks+building pages}

每当 \LUATEX\ 准备好将内容移至主垂直列表时，就会调用此回调。你可以
使用此回调对页面构建阶段进行特殊处理，例如拼版或栏平衡。

\begin{typing}
function(<string> extrainfo)
end
\end{typing}

字符串 \type {extrainfo} 提供有关 \TEX 相对于 ``当前页面 的状态的
一些额外信息。\cbk {buildpage_filter} 回调的可能值为：

\begin{longtable}{l p{5cm}}
\toprule 值                   & 说明                              \\
\midrule
\type{alignment}       & 正在添加一个（部分）对齐    \\
\type{after_output}    & 输出例程刚刚结束             \\
\type{new_graf}        & 新段落的开始                 \\
\type{vmode_par}       & 在垂直模式下遇到了 \prm {par} \\
\type{hmode_par}       & 在水平模式下遇到了 \prm {par} \\
\type{insert}          & 添加了一个插入               \\
\type{penalty}         & 一个处罚点（在垂直模式下）   \\
\type{before_display}  & 就在显示公式开始之前         \\
\type{after_display}   & 显示公式结束                 \\
\type{end}             & \LUATEX\ 正在终止（一切都结束了） \\
\bottomrule
\end{longtable}

\subsection{\cbk {build_page_insert}}

\topicindex{callbacks+inserts}

当页面构建器添加插入时，会调用此回调。对这个机制没有太多控制权，
但此回调允许在插入之前对间距进行一些最后一刻的调整，例如当连续追加
多个插入（类型）时，这可能会很方便。

\begin{typing}
function(<number> n, <number> i)
    return <number> register
end
\end{typing}

其中

\begin{longtable}{l p{5cm}}
\toprule 值     & 说明              \\
\midrule
\type{n} & 插入类别         \\
\type{i} & 插入的顺序       \\
\bottomrule
\end{longtable}

返回值是一个数字，指示用于前置间距的 skip 寄存器。例如，这允许使用
不同的顶部间距（当 \type {i} 等于一时）和中间间距（当 \type {i} 大于一时）。
当然，你可以随意处理插入盒子，但你需要确保 \LUATEX\ 之后能正常工作。

\subsection{\cbk {pre_linebreak_filter}}

\topicindex{callbacks+linebreaks}

此回调在 \LUATEX\ 开始将节点列表转换为 \prm {hbox} 堆栈之前、
即添加 \prm {parfillskip} 之后被调用。

\begin{typing}
function(<node> head, <string> groupcode)
    return true | false | <node> newhead
end
\end{typing}

名为 \type {groupcode} 的字符串标识节点列表在 \TEX 处理过程中的上下文。
可能的取值范围在下表中给出，但并非所有这些值都实际出现在
\cbk {pre_linebreak_filter} 中，其中一些属于 \cbk {hpack_filter} 和
\cbk {vpack_filter} 回调，将在接下来的两段中解释。

\begin{longtable}{l p{5cm}}
\toprule 值                 & 说明                                   \\
\midrule
\type{<empty>}       & 主垂直列表                           \\
\type{hbox}          & 水平模式下的 \prm {hbox}            \\
\type{adjusted_hbox} & 垂直模式下的 \prm {hbox}             \\
\type{vbox}          & \prm {vbox}                            \\
\type{vtop}          & \prm {vtop}                            \\
\type{align}         & \prm {halign} 或 \prm {valign}          \\
\type{disc}          & 断点                                  \\
\type{insert}        & 打包一个插入                          \\
\type{vcenter}       & \prm {vcenter}                         \\
\type{local_box}     & \lpr {localleftbox} 或 \lpr {localrightbox} \\
\type{split_off}     & \prm {vsplit} 的顶部                   \\
\type{split_keep}    & \prm {vsplit} 的剩余部分               \\
\type{align_set}     & 对齐单元                              \\
\type{fin_row}       & 对齐行                                \\
\bottomrule
\end{longtable}

与所有处理节点的回调一样，返回值可以是以下三种之一：

\begin{itemize}\item 布尔值 \type {true} 表示处理成功

\item \type {<node>} 表示应将 ``head 节点替换为返回的节点

\item 布尔值 \type {false} 表示应忽略 ``head 节点列表并将其从
    内存中清除

\end{itemize}

此回调不替换任何内部代码。

\subsection{\cbk {linebreak_filter}}

\topicindex{callbacks+linebreaks}

此回调替换 \LUATEX 的断行算法。

\begin{typing}
function(<node> head, <boolean> is_display)
    return <node> newhead
end
\end{typing}

返回的节点是将被添加到主垂直列表的列表头；如果此段落被后续的数学
显示公式打断，则布尔参数为 true。

如果你返回的不是 \type {<node>}，\LUATEX\ 将对从 \type {<head>} 开始的
列表应用内部断行算法。否则，你返回的 \type {<node>} 应该是一个节点列表
的表头，这些节点都在垂直模式下允许，并且其中至少有一个必须表示一个
hbox。否则将导致致命错误。

将此回调设置为 \type {false} 是可能的，但很危险，因为你可能最终陷入
一个无法修复的 ``死循环（deadcycles loop）。

\subsection{\type {append_to_vlist_filter}}

\topicindex{callbacks+contributions}

每当 \LUATEX\ 向垂直列表添加盒子时，就会调用此回调：

\begin{typing}
function (<node> box, <string> locationcode, <number> prevdepth, <boolean> mirrored)
    return list, prevdepth
end
\end{typing}

不返回任何值也是可以的，在这种情况下你还需要清除该盒子或自行处理它。
prevdepth 也是可选的。位置有 \type {box}、\type {alignment}、
\type {equation}、\type {equation_number} 和 \type {post_linebreak}。
你可以传递 \type {nil} 而不是节点。

\subsection{\cbk {post_linebreak_filter}}

\topicindex{callbacks+linebreaks}

此回调在 \LUATEX\ 将节点列表转换为 \prm {hbox} 堆栈之后立即被调用。

\begin{typing}
function(<node> head, <string> groupcode)
    return true | false | <node> newhead
end
\end{typing}

此回调不替换任何内部代码。

\subsection{\cbk {hpack_filter}}

\topicindex{callbacks+packing}

当 \TEX 准备好开始将某些水平模式材料打包成盒子时，会调用此回调。
目前数学项和行盒子被忽略。

\begin{typing}
function(<node> head, <string> groupcode, <number> size,
         <string> packtype [, <string> direction] [, <node> attributelist])
    return true | false | <node> newhead
end
\end{typing}

\type {packtype} 为 \type {additional} 或 \type {exactly}。如果是
\type {additional}，则 \type {size} 是一个 \type {\hbox spread ...} 参数。
如果是 \type {exactly}，则 \type {size} 是一个 \type {\hbox to ...}。
在两种情况下，该数字都以缩放点为单位。

\type {direction} 是三个字母的方向指示符字符串之一，或者是 \type {nil}。

此回调不替换任何内部代码。

\subsection{\cbk {vpack_filter}}

\topicindex{callbacks+packing}

当 \TEX 准备好开始将某些垂直模式材料打包成盒子时，会调用此回调。
目前数学显示公式被忽略。

此函数与 \cbk {hpack_filter} 非常相似。除了调用时机不同之外，还有一个
额外的变量，对应于 \TEX 的 \prm {maxdepth} 设置。

\begin{typing}
function(<node> head, <string> groupcode, <number> size, <string> packtype,
        <number> maxdepth [, <string> direction] [, <node> attributelist])
    return true | false | <node> newhead
end
\end{typing}

此回调不替换任何内部代码。

\subsection{\type {hpack_quality}}

\topicindex{callbacks+packing}

此回调可用于拦截因打包水平列表（如段落构建器中发生的）而可能产生的
overfull 消息。该函数接受几个参数：

\begin{typing}
function(<string> incident, <number> detail, <node> head, <number> first,
         <number> last)
    return <node> whatever
end
\end{typing}

incident 是 \type {overfull}、\type {underfull}、\type {loose} 或
\type {tight} 之一。detail 要么是 \type {overfull} 时的溢出量，要么是
其他情况下的坏度（badness）。head 是所构建的列表（当启用突出或扩展时，
这是一个中间列表）。你可以选择返回一个节点，例如 overfull 规则指示器。
该节点将被追加到列表中（就像 \TEX 自己的规则一样）。

\subsection{\type {vpack_quality}}

\topicindex{callbacks+packing}

此回调可用于拦截因打包垂直列表（如页面构建器中发生的）而可能产生的
overfull 消息。该函数接受几个参数：

\begin{typing}
function(<string> incident, <number> detail, <node> head, <number> first,
         <number> last)
end
\end{typing}

incident 是 \type {overfull}、\type {underfull}、\type {loose} 或
\type {tight} 之一。detail 要么是 \type {overfull} 时的溢出量，要么是
其他情况下的坏度。head 是所构建的列表。

\subsection{\cbk {process_rule}}

\topicindex{callbacks+rules}

这是一个实验性回调。它可用于子类型~4（用户）的规则。该回调接收三个
参数：节点、宽度和高度。该回调可以使用 \type {pdf.print} 向 \PDF\ 文件
写入代码，但请注意不要破坏最终结果。不进行任何检查。

\subsection{\type {pre_output_filter}}

\topicindex{callbacks+output}

当 \TEX 准备好开始为 \prm {output} 打包盒子 255 时，会调用此回调。

\begin{typing}
function(<node> head, <string> groupcode, <number> size, <string> packtype,
        <number> maxdepth [, <string> direction])
    return true | false | <node> newhead
end
\end{typing}

此回调不替换任何内部代码。

\subsection{\cbk {hyphenate}}

\topicindex{callbacks+hyphenation}

\begin{typing}
function(<node> head, <node> tail)
end
\end{typing}

无返回值。此回调必须在其接收到的节点列表中插入可选断点节点。

将此回调设置为 \type {false} 将阻止内部的可选断点插入过程。

\subsection{\cbk {ligaturing}}

\topicindex{callbacks+ligature building}

\begin{typing}
function(<node> head, <node> tail)
end
\end{typing}

无返回值。此回调必须对其接收到的节点列表应用连字处理。

你不必担心返回值，因为传递给回调的 \type {head} 节点保证不是 glyph_node
（如有需要，会前置一个临时节点），因此它不会受到所发生变更的影响。
回调之后，``列表尾部 的内部值将被重新计算。

\type {head} 的 \type {next} 保证非 nil。

\type {tail} 的 \type {next} 保证为 nil，因此第二个回调参数通常可以忽略。
提供它是为了保持一致性，并且在需要进行特殊处理时有时会很方便。

将此回调设置为 \type {false} 将阻止内部的连字创建过程。

你不得破坏节点列表。例如，head 通常是一个 local par 节点，而 tail 是
一个 glue。过度破坏可能会使 \LUATEX\ 进入恐慌模式。

\subsection{\cbk {kerning}}

\topicindex{callbacks+kerning}

\begin{typing}
function(<node> head, <node> tail)
end
\end{typing}

无返回值。此回调必须在其接收到的节点列表中的节点之间应用字距调整。
调用约定参见 \cbk {ligaturing}。

将此回调设置为 \type {false} 将阻止内部的 kern 插入过程。

你不得破坏节点列表。例如，head 通常是一个 local par 节点，而 tail 是
一个 glue。过度破坏可能会使 \LUATEX\ 进入恐慌模式。

\subsection{\type {insert_local_par}}

每个段落都以一个 local par 节点开始，该节点用于跟踪例如方向等信息。
你可以在创建过程中挂接一个回调：

\begin{typing}
function(<node> local_par, <string> location)
end
\end{typing}

没有返回值，你应该确保节点保持有效，否则 \TEX 可能会混乱。

\subsection{\cbk {mlist_to_hlist}}

\topicindex{callbacks+math}

此回调替换 \LUATEX 的数学列表到节点列表的转换算法。

\begin{typing}
function(<node> head, <string> display_type, <boolean> need_penalties)
    return <node> newhead
end
\end{typing}

返回的节点是将被添加到垂直或水平列表的列表头；字符串参数根据当前
数学模式为 ``text 或 ``display；如果需要在此列表中插入
处罚点，布尔参数为 \type {true}，否则为 \type {false}。

将此回调设置为 \type {false} 很糟糕，几乎肯定会导致死循环。

\section{信息报告回调}

\subsection{\cbk {pre_dump}}

\topicindex{callbacks+dump}

\begin{typing}
function()
end
\end{typing}

此函数在开始转储到格式文件之前被调用。它不替换任何代码，既没有参数
也没有返回值。

\subsection{\cbk {start_run}}

\topicindex{callbacks+job run}

\begin{typing}
function()
end
\end{typing}

此回调替换打印 \LUATEX 横幅的代码。请注意，要成功使用，此回调必须在
\LUA\ 初始化脚本中设置，否则它只会在运行已经开始后才被看到。

\subsection{\cbk {stop_run}}

\topicindex{callbacks+job run}

\begin{typing}
function()
end
\end{typing}

此回调替换打印 \LUATEX 统计信息和 ``output written to 消息的代码。
引擎仍可能进行内部整理，因此你不应依赖此挂钩来后处理 \PDF\ 或日志文件。

\subsection{\cbk {start_page_number}}

\topicindex{callbacks+pages}

\begin{typing}
function()
end
\end{typing}

替换在 \prm {shipout} 开始时打印 \type {[} 和页码的代码。此回调还将
覆盖通常在 \prm {tracingoutput} 为正时进行的盒子信息打印。

\subsection{\cbk {stop_page_number}}

\topicindex{callbacks+pages}

\begin{typing}
function()
end
\end{typing}

替换在 \prm {shipout} 结束时打印 \type {]} 的代码。

\subsection{\cbk {show_error_hook}}

\topicindex{callbacks+errors}

\begin{typing}
function()
end
\end{typing}

此回调从 \TEX 错误函数内部运行，其目的是允许你在 \TEX 已有的操作之上
进行一些额外的报告（不会移除任何正常操作）。你可能会发现 \type {status}
表中的一些值很有用。此回调不替换任何内部代码。

\subsection{\cbk {show_error_message}}

\topicindex{callbacks+errors}

\begin{typing}
function()
end
\end{typing}

此回调替换打印错误消息的代码。消息之后的常规交互不受影响。

\subsection{\cbk {show_ignored_error_message}}

\topicindex{callbacks+errors}

\begin{typing}
function()
end
\end{typing}

此回调替换在启用 \prm {ignoreprimitiveerror} 时打印错误消息的代码。
与之前一样，消息之后的常规交互不受影响。

\subsection{\cbk {show_lua_error_hook}}

\topicindex{callbacks+errors}

\begin{typing}
function()
end
\end{typing}

此回调替换打印额外 \LUA\ 错误消息的代码。

\subsection{\cbk {start_file}}

\topicindex{callbacks+files}

\begin{typing}
function(category,filename)
end
\end{typing}

此回调替换在打开文件时打印 \LUATEX 标记的代码，例如对于普通文件打印
\type {(filename}。category 是一个数字：

\begin{longtable}{c l}
\toprule 值   & 含义 \\
\midrule
1 & 普通数据文件，如 \TEX\ 源文件 \\
2 & 字体映射，将字体名称关联到资源 \\
3 & 图像文件（\type {png}、\type {pdf} 等） \\
4 & 嵌入的字体子集 \\
5 & 完全嵌入的字体 \\
\bottomrule
\end{longtable}

\subsection{\cbk {stop_file}}

\topicindex{callbacks+files}

\begin{typing}
function(category)
end
\end{typing}

此回调替换在关闭文件时打印 \LUATEX 标记的代码，例如对于普通文件打印
\type {)}。

\subsection{\cbk {call_edit}}

\topicindex{callbacks+editing}

\begin{typing}
function(filename,linenumber)
end
\end{typing}

当在错误消息提示下按下 ``E 时，此回调替换对外部编辑器的调用。
回调将控制权返回给主程序后，处理将立即结束。

\subsection{\cbk {finish_synctex}}

\topicindex{callbacks+synctex}

此回调可用于完成替代的 synctex 方法。它在正常的 synctex 终结器之后
启动（当未启用原生 synctex 时，该终结器会在一次运行后删除 synctex 文件）。

\subsection{\cbk {wrapup_run}}

\topicindex{callbacks+wrapping up}

此回调在 \PDF\ 和日志文件关闭后被调用。使用风险自负。

\section{与 \PDF\ 相关的回调}

\subsection{\cbk {finish_pdffile}}

\topicindex{callbacks+\PDF\ file}

\begin{typing}
function()
end
\end{typing}

当所有文档页面已写入 \PDF\ 文件，且 \LUATEX\ 即将完成输出文档结构时，
会调用此回调。其预期用途是最终更新 \PDF\ 字典，例如 \type {/Catalog}
或 \type {/Info}。此回调不替换任何代码。既没有参数也没有返回值。

\subsection{\cbk {finish_pdfpage}}

\topicindex{callbacks+\PDF\ file}

\begin{typing}
function(shippingout)
end
\end{typing}

此回调在 \PDF\ 页面流组装完成之后、页面对象最终确定之前被调用。

\subsection{\cbk {page_order_index}}

\topicindex{callbacks+\PDF\ file}

这是供专家使用的一个回调，用来操纵页面树——一种决定 \PDF\ 文件中
页面顺序的数据结构：

\begin{typing}
function(pagenumber)
  return pagenumber
end
\end{typing}

假设我们有 12 页，那么我们可以这样做：

\begin{typing}
callback.register("page_order_index",function(page)
        if page ==  1 then return 12
    elseif page ==  2 then return 11
    elseif page == 11 then return  2
    elseif page == 12 then return  1
    else                   return page
    end
end)
\end{typing}

这将交换前两页和后两页。你需要知道页数，这是实现的一个副作用。
当你搞砸了的时候，不要抱怨。

\subsection{\cbk {process_pdf_image_content}}

\topicindex{callbacks+image content}

当嵌入 \PDF\ 文件中的一页时，其页面流以及相关对象会被复制到目标文件中。
但是，页面流中可能包含假定有额外资源的运算符，例如标记文本。你可以
决定对其进行过滤，为此 \LUATEX\ 提供了一个回调。下面是一个简单的
使用示例：

\begin{typing}
pdf.setrecompress(1)

callback.register("process_pdf_image_content",function(s)
    print(s)
    return s
end)
\end{typing}

你需要显式启用重新压缩，否则内容流将仅以其原始压缩形式传递。

\section{与字体相关的回调}

\subsection{\cbk {define_font}}

\topicindex{callbacks+fonts}

\begin{typing}
function(<string> name, <number> size, <number> id)
    return <table> font | <number> id
end
\end{typing}

字符串 \type {name} 是字体规范中的文件名部分，由用户给出。

数字 \type {size} 有点特殊：

\begin{itemize}
\item 如果为正数，则以缩放点指定一个 ``at size（指定字号）。

\item 如果为负数，其绝对值表示相对于字体设计字号的 ``scaled
    （缩放）设置。

\end{itemize}

\type {id} 是分配给字体的内部编号。

要返回的 \type {font} 表的内部结构在 \ref{fonts} 中解释。
该表在内部保存，因此你可以在表中放置额外字段供你后续的 \LUA\ 代码使用。
或者，\type {retval} 可以是先前定义的 fontid。如果可以重用先前的定义
而不是创建全新的字体结构，这会很有用。

将此回调设置为 \type {false} 没有意义，因为它会完全阻止字体加载，
但仍然会生成错误。

\subsection{\cbk {glyph_not_found} 与 \cbk {glyph_info}}

\topicindex{callbacks+fonts}

\type {glyph_not_found} 回调在设置后，会在后端无法插入字形时启动。
如果没有定义回调，则会向日志写入一条消息。

\begin{typing}
function(<number> id, <number> char)
    -- do something with font id and char code
end
\end{typing}

可以设置 \type {glyph_info} 回调来报告字形的有用表示。

\begin{typing}
function(<node> g)
    -- return a string or nil
end
\end{typing}

当返回 \type {nil} 时，打印字符代码；否则使用返回的字符串。默认显示
\UTF\ 表示，这不总是那么有用，尤其是在没有实际表示时。请记住，设置
此回调可能会以不兼容的方式更改日志。

\subsection{\cbk {provide_charproc_data}}

\topicindex{callbacks+fonts}

当后端正在写出用户定义的 Type~3 字体时，会触发 \type
{provide_charproc_data} 回调：

\begin{typing}
function(<number> mode, <number> id, <number> char)
    if mode == 1 then
        -- preroll
        return
    elseif mode == 2 then
        -- write out the glyph contents
        return <number> stream, <number> width
    elseif mode == 3 then
        -- the overall font scale
        return <number> scale
    end
end
\end{typing}

要创建用户定义的 Type~3 字体，你需要在定义字体时在顶层设置
\type {encodingbytes = 0} 和 \type {psname = "none"}。从
\type {glyph_not_found} 回调中，你应在 mode~3 时返回整体字体缩放比例
（通常为 0.001），在 mode~2 时返回 \PDF\ 流的索引（其中第一个运算符为
\type {d0} 或 \type {d1}）以及字形的宽度（以 sp 为单位）。
通常可以忽略 mode~1。
